\cleardoublepage
\thispagestyle{plain}
\addcontentsline{toc}{chapter}{SUMMARY}
\begingroup
\singlespacing

\vspace*{72pt}
\begin{center}
  \bfseries IMPROVED T-TYPE CONVERTER ARCHITECTURE FOR ENERGY STORAGE SYSTEMS WITH SWITCHING AND SIN\textsuperscript{3}-BASED COMPENSATION METHODS IMPROVING LOW-POWER PERFORMANCE\\[18pt]
  SUMMARY
\end{center}

Energy storage systems are among the key enablers of renewable integration. Their overall performance is determined not only by the battery pack but also by the efficiency, reliability, volume and cost of the power conversion interface between the battery, the DC bus and the grid. Grid-connected photovoltaic and storage converters operate at rated power only for limited periods and spend a considerable share of their operating life at partial load, where fixed driving, switching and magnetic losses grow relative to the transferred power and the current waveform becomes more sensitive to switch nonlinearities, dead time and grid distortion. The low-power region is therefore the operating range in which the trade-off between efficiency and current total harmonic distortion (THD) becomes most visible.

The three-level T-type inverter is an attractive grid interface owing to its high efficiency and power density at rated load. In the low-power region, however, its behaviour is governed by a different mechanism. The bidirectional neutral branch that produces the zero-voltage state keeps conducting after the inductor current reaches zero, so the phase leg behaves like a synchronous buck converter and lets the current reverse within a switching period. The reverse circulating current carries no useful power, yet it increases the RMS current and the conduction losses of the semiconductors and filter elements. Blocking reverse current is a known principle for synchronous buck converters, but applying it to the zero state of a grid-connected three-level T-type inverter changes the current path, the pole voltage, the modulation integrity and the harmonic content at the same time, and is therefore a topology-specific problem.

The thesis addresses the problem with three original contributions built on the T-type topology, which correspond to the three components of the title. The first two operate at the phase-leg scale, the third at the system scale. (1) A neutral-switch suppression method (the \emph{switching method improving low-power performance}) blocks the neutral-switch command that sustains reverse conduction. (2) A Sin\textsuperscript{3}-based feedforward compensation (the \emph{Sin\textsuperscript{3}-based compensation}) reduces the third-harmonic penalty of the suppression without altering the stability of the current loop. (3) The GT-type reconfigurable energy storage architecture (the \emph{improved T-type converter architecture for energy storage systems}) re-tasks the same T-type semiconductors as a grid inverter, a Vienna-type power factor correction (PFC) stage and an active battery management stage. Four research questions are answered: the mechanism of reverse circulation, the efficiency--harmonic trade-off of suppression, the harmonic effect and loop interaction of the Sin\textsuperscript{3} term, and the feasibility of safe hardware re-tasking.

\textbf{System, model and method.} The prototype is a three-phase, four-wire T-type inverter whose neutral is tied to the DC-bus midpoint; simulations and experiments are carried out on one phase leg with an LCL filter ($L_1=600~\mu$H, $C_f=8~\mu$F, grid-side equivalent $L_2$), a switching frequency of 10~kHz and a dead time of 1~$\mu$s read from the PWM timer registers. The switching states, the four basic current paths and the synchronous-buck equivalent of the phase leg are drawn on a common circuit base, and the reverse-current threshold is derived from the condition that the average inductor current is smaller than half of the ripple. The evidence comes from four switched Simscape models: classical complementary switching (S1), neutral-switch suppression (S2), suppression with Sin\textsuperscript{3} feedforward (S3), and a closed-loop current-feedback model (S4) built in the thesis on top of S3. The gate logic read from the models confirms that the only difference between S1 and S2 is the opposite neutral-switch command in the zero state. Before use, the models were audited against the thesis definition of the circuit; the problems found, such as an EMC capacitor not matching its single-phase equivalent and mis-mapped diodes, were corrected with the smallest change consistent with the topology, and the effect of every correction was measured. All circuit, semiconductor, modulation, control and solver parameters are collected in a single table with value, unit, source and type (measured, datasheet, design, calibrated, simulation choice, corrected).

\textbf{Calibration against the experiment.} The S1--S3 models were calibrated to the experimental summary with only two free variables per scenario, the modulation amplitude and the reference phase. The targets were the active power (within $\pm1$\,\%) and the magnitude of the fundamental reactive power (within the rounding band of the analyser's $\cos\varphi$ reading). Circuit parameters were not changed; AC current, efficiency, DC power and THD were kept as independent checks. Calibration showed that the bridge does not behave as an ideal voltage source at low current: its Thevenin impedance, measured from the runs, is 1.67~$\Omega$ in S1 because of dead time and 28.7~$\Omega$ and 13.7~$\Omega$ in S2 and S3 because of discontinuous conduction, which occupies 12.4\,\% and 22.6\,\% of the window. As a consequence, the active power depends mainly on the amplitude and the reactive power mainly on the phase, and the two targets were solved jointly with an empirical Jacobian.

\textbf{Low-power results.} At a phase current of about 2.3~A, 250--280~W and a grid voltage of about 120~V, suppression raises the efficiency from 92.5\,\% to 94.2\,\% in the experiment and from 91.15\,\% to 93.93\,\% in the calibrated simulation. The ranking is the same in both sources, and the larger model gain stems from a $-1.35$ point deviation of the S1 efficiency. The price is a pronounced third-harmonic distortion: the grid-current THD reaches 47\,\% in the experiment and rises from 39.5\,\% to 56.7\,\% in the calibrated model. With the Sin\textsuperscript{3} term added, the experimental THD falls to 14\,\% and the efficiency reaches 94.5\,\%. The calibrated S3 model reproduces the AC current within 1.4\,\%, the efficiency within 0.25 points and, without THD being a calibration target, the current THD as 14.8\,\%, which is the strongest independent evidence that the model captures the low-current distortion with the correct magnitude. The reported S1 experimental THD of 15\,\% is inconsistent with the reported power factor of 0.922, which requires about 42\,\% THD at unity displacement factor; the model agrees with the power, voltage and current of the S1 record. In S2 the model overestimates the AC current by 3.7\,\% and the THD by 9.7 points, and the source of this gap is reported as a limitation.

\textbf{Sin\textsuperscript{3} compensation.} In the model the term is the cube of a unit-amplitude sine, added with positive sign and without normalisation to the modulation sum as an open-loop feedforward, so that $m(\theta)=K_1\sin\theta+K_3\sin^3\theta$. Using the identity $\sin^3\theta=[3\sin\theta-\sin(3\theta)]/4$, the coefficients of $\sin^3\theta$ and of a direct $\sin(3\theta)$ injection are shown to carry different meanings, and simplified derivations are shown to lie outside their validity range because a single operating-point ratio cannot be treated as a transfer function. A piecewise DCM--CCM average model explains why a sinusoidal duty cycle does not produce a sinusoidal current after suppression. The scan with the correct identity reduces the analytical third-harmonic ratio from 39.96\,\% for a pure sinusoidal duty cycle to 0.15\,\% at $(D_1,D_3)=(0.34,0.40)$; the analytical third-harmonic coefficient ($-0.10$) is close to the calibrated coefficient of the switched model ($-0.096$), and the coefficients depend on the operating point.

\textbf{Closed loop at medium and rated current.} The S4 model adds L1-current feedback with P, PI and PIR controllers in an analogue-equivalent form and an enlarged DC bus. The enlargement is needed because the original bus capacitors would let the midpoint swing by about 120~V at 25~A; the enlarged bus reduces the swing to about 10~V. At the rated 25~A, PI and PIR controllers track the reference with $-2.1$\,\% and $-1.0$\,\% amplitude error and 8.5--9.2\,\% THD without saturation, while the P controller leaves $-11.4$\,\%. For a 10~A to 25~A reference step the overshoot relative to the final steady envelope is 0.2\,\% and the settling time 0.44~ms. At 10~A, however, the suppressed bridge is in discontinuous conduction over most of the cycle, the THD is 19.6--22.6\,\% and the P and PI controllers leave $-23$\,\% and $-11$\,\% amplitude error. Suppression is therefore a low-power method. At 25~A the Sin\textsuperscript{3} gain calibrated at 2.3~A increases the third harmonic, so the feedforward amplitude must be scheduled with current.

\textbf{Stability.} The stability analysis shows in three layers that the Sin\textsuperscript{3} term is a known input added from outside the loop and not a controller. Analytically, the characteristic equation $1+L(s)=0$ does not contain the term. In the averaged model and in the linearisation of the switched model, the loop gain $L$, the sensitivity $S$ and the complementary sensitivity $T$ are identical with the term on and off, and in the switched model the differences remain at the level of measurement uncertainty. The effect of the term on the current is set by $S(j\omega)K_{\mathrm{pwm}}G(j\omega)$, and its 150~Hz component is attenuated by $|S|=0.23$. Because the loop gain crosses unity three times (586, 2454 and 3386~Hz), margins are given at every crossing and the verdict rests on the Nyquist criterion and the closed-loop poles. The analogue-equivalent loop at 25~A with PI control has a critical phase margin of 54° at 3.4~kHz, a gain margin of +35~dB at 22~kHz, a delay margin of 44~$\mu$s and a sensitivity peak of 2.4~dB. Its least-damped pole lies at 3.1~kHz with $\zeta=0.18$, and the loop remains stable for 0--5~mH of additional grid inductance. A digital implementation that samples at the start of each PWM period and applies the result one period later is unstable in this LCL filter, whose resonance (2.8~kHz) lies above $f_s/6$: the dominant pole has $|z|=1.18$, and double update still gives $|z|=1.085$. A hardware implementation therefore needs active damping or grid-current feedback.

\textbf{GT-type architecture.} The third contribution re-tasks the T-type switches according to the operating mode, combining the inverter, Vienna PFC and active battery management functions on a common hardware base. A first design layer shows the function-sharing paths and the charging chain in simulation (PFC, DC bus, transformer primary, secondary and buck/battery). A 100~kW project layer extends the concept with AB--BC--CA battery groups, isolated per-battery charging branches and a selectable series discharge path. Operating modes are defined as safe off, pre-charge, PFC/CC--CV charging, staged discharge to load, bypass/backup and a chemistry-dependent maintenance mode. Safe reconfiguration requires near-zero current, node-voltage matching, dead time, contactor sequencing and switch feedback at the same time. Cost, volume and loss advantages over separate power stages are stated as potential, because no prototype, bill of materials or thermal model exists for the target design.

\textbf{Conclusions and limitations.} The thesis shows that the neutral-branch behaviour of the T-type converter at low power offers an efficiency gain with an explicit harmonic trade-off, that the Sin\textsuperscript{3} feedforward reduces the harmonic penalty without changing the stability of the current loop, and that the same T-type semiconductor base can be organised into a reconfigurable storage architecture. Simulation and experimental results are reported separately and the evidence level of every claim is stated. The main limitations are the following. Experiments exist only at about 2.3~A, as screen records and a one-line summary per method without raw data or repetitions, and no new experiment was performed. The medium- and rated-current results rely on simulation with an enlarged DC bus and analogue-equivalent control. $L_2$ is an equivalent of the grid and variac leakage that was not measured directly, and $L_1$ saturation is not modelled. The switching-mode transition from suppression to complementary switching at higher current, current-dependent scheduling of the feedforward amplitude, a digital controller with active damping and a GT-type prototype are proposed as future work.

\textbf{Keywords:} T-type inverter, energy storage system, low-power operation, neutral-switch suppression, feedforward compensation, third harmonic, stability analysis, reconfigurable architecture.
\endgroup
